% 原书分章PDF第16—19页：第26.4节导言与第26.4.1节。
\section{示例：g计算参数的估计}

我们通过构造2-TMLE，给出上述算法的一个具体例子。根据 $O=(W,A,Y)\sim P_0$ 的 $n$ 个独立副本，估计 $\psi_0:=\Psi(P_0)$，其中 $\Psi(P)=E_PE_P(Y\mid A=1,W)$ 为单时点干预的g计算参数。在不可检验的因果假设下，这一参数对应于单时点干预 $A=1$ 下反事实结局的均值。

以 $o:=(w,a,y)$ 表示 $O$ 的一个可能实现值，记 $\bar g(w):=P(A=1\mid W=w)$、$\bar Q(w):=E_P(Y\mid A=1,W=w)$，并以 $Q_W$ 表示 $W$ 边际分布关于适当计数测度的密度。已知此时的一阶典范梯度为
\[
P\mapsto D^{(1)*}(P):=D_Y^{(1)*}(P)+D_W^{(1)*}(P),
\]
其中逐点定义
\[
\begin{aligned}
D_Y^{(1)*}(P)(o)&:=H^{(1)}(\bar g)(a,w)[y-\bar Q(w)],\\
D_W^{(1)*}(P)(o)&:=\bar Q(w)-\Psi(P),
\end{aligned}
\]
且 $H^{(1)}(\bar g)(a,w):=a/\bar g(w)$、$\bar g(w):=P(A=1\mid W=w)$。根据这一定义，可得 $\Psi(P)-\Psi(P_0)=-P_0D^{(1)*}(P)+R_2(P,P_0)$，其中
\[
R_2(P,P_0):=P_0\left[\left(\frac{\bar g-\bar g_0}{\bar g}\right)(\bar Q-\bar Q_0)\right],
\]
$\bar g_0$ 和 $\bar Q_0$ 分别表示 $P_0$ 下的 $\bar g$ 和 $\bar Q$。我们的目标是找到表示 $R_2(P,P_0)=-\frac12P_0^2D^{(2)}(P)+R_3(P,P_0)$。下面分别处理 $W$ 具有有限支撑集和无限支撑集的情形。

\subsection{情形I：有限支撑集}

假设 $W$ 在每个所考虑的 $P$ 下均具有有限支撑集。可以证明，$P$ 处、在 $(o_1,o_2)$ 点的二阶典范梯度，由下式的对称化给出，其中 $o_1:=(w_1,a_1,y_1)$、$o_2:=(w_2,a_2,y_2)$：
\begin{equation}
D^{(2)}(P)(o_1,o_2):=H^{(2)}(P)(w_1,a_1,w_2,a_2)[y_1-\bar Q(w_1)],
\tag{26.6}
\end{equation}
这里
\[
H^{(2)}(P)(w_1,a_1,w_2,a_2):=
\frac{2a_1I(w_1=w_2)}{\bar g(w_1)Q_W(w_1)}
\left[1-\frac{a_2}{\bar g(w_1)}\right].
\]
不难直接验证，确实有
\begin{equation}
\frac12P_0^2D^{(2)}(P)=
-P_0\left[\left(\frac{\bar g-\bar g_0}{\bar g}\right)(\bar Q-\bar Q_0)\right]+R_3(P,P_0),
\tag{26.7}
\end{equation}
其中 $R_3(P,P_0)$ 为
\[
P_0\left[\left(1-\frac{\bar g_0Q_{W,0}}{\bar gQ_W}\right)
\left(\frac{\bar g-\bar g_0}{\bar g}\right)(\bar Q-\bar Q_0)\right].
\]
因此，对于一个三阶项 $R_3(P,P_0)$，如我们所需，有
% 原书这里实际印P_0D^(2)，没有平方；只改行内为无编号展示。
\[
\Psi(P)-\Psi(P_0)=-P_0D^{(1)*}(P)-\frac12P_0D^{(2)}(P)+R_3(P,P_0).
\]
由于 $W$ 的支撑集无限，事件 $\{W_1=W_2\}$ 以正概率发生，因而 $D^{(2)}(P)$ 不退化。

在这一情形下，目标参数 $\Psi(P)$ 通过 $Q(P)=(\bar Q(P),Q_W(P))$ 依赖于 $P$。不难验证，损失函数 $L(Q):=L_1(\bar Q)+L_2(Q_W)$ 使 $(\bar Q_0,Q_{W,0}):=(\bar Q(P_0),Q_W(P_0))$ 确实达到真实风险 $P_0L(Q)$ 的最小值，其中对每个 $o=(w,a,y)$，逐点定义
\[
\begin{aligned}
L_1(\bar Q)(o)&:=-a[y\log\bar Q(w)+(1-y)\log(1-\bar Q(w))],\\
L_2(Q_W)(o)&:=-\log Q_W(w).
\end{aligned}
\]

令 $g(P):=\bar g$，则 $D^{(1)*}(P)$ 和 $D^{(2)}(P)$ 都通过 $(Q(P),g(P))$ 依赖于 $P$。由于 $(Q(P),g(P))$ 的各个分量涉及似然中彼此正交的部分，它们相互变动独立。这为构造适当的扰动子模型提供了更大灵活性。定义
\[
\bar D_n^{(2)}(P)(o):=\bar H_n^{(2)}(P)(w,a)[y-\bar Q(w)],
\]
其中 $\bar H_n^{(2)}(P)(w,a):=n^{-1}\sum_{j=1}^nH^{(2)}(P)(w,a,W_j,A_j)$。注意，$\bar D_n^{(2)}(P)\in T_{\mathcal M}(P)$，而且它实际上是给定 $A$ 和 $W$ 时 $Y$ 的条件分布的一个得分。此外，
\[
P_n\bar D_n^{(2)}(P)=P_n^2D^{(2)}(P).
\]
还可注意到，在 $Q_W=Q_{W,n}$，即 $W$ 的经验分布处，
\[
\bar H_n^{(2)}(P)(w,a)=\frac{2a}{\bar g(w)}
\left[1-\frac{\bar g_{n,\mathrm{NP}}(w)}{\bar g(w)}\right],
\]
其中 $\bar g_{n,\mathrm{NP}}(w):=\sum_{i=1}^nI(A_i=1,W_i=w)/\sum_{i=1}^nI(W_i=w)$ 是 $\bar g_0(w)$ 的非参数极大似然估计量。

容易验证，若 $Q_W(P)=Q_{W,n}$，则 $P_nD_W^{(1)*}(P)=0$。因此，如果取 $Q_{W,n}$ 为 $Q_{W,0}$ 的初始估计量，要达到求解所需一阶与二阶估计方程的目标，只需给出 $\bar Q_0$ 的估计 $\bar Q_n^*$，使其满足
\begin{equation}
P_nD_Y^{(1)*}(\bar Q_n^*,Q_{W,n},\bar g_n)
=P_n\bar D_n^{(2)}(\bar Q_n^*,Q_{W,n},\bar g_n)=0,
\tag{26.8}
\end{equation}
其中 $\bar g_n$ 是我们对 $\bar g_0$ 的初始估计。给定任意特定的 $\bar Q$、$Q_W$ 和 $\bar g$，由
\[
\bar Q_{\bar g,Q_W}(\epsilon):=\expit\left[\logit(\bar Q)
+\epsilon_1H^{(1)}(\bar g)+\epsilon_2\bar H_n^{(2)}(\bar g,Q_W)\right]
\]
确定的子模型，令 $\epsilon:=(\epsilon_1,\epsilon_2)$，满足 $\bar Q_{\bar g,Q_W}(0)=\bar Q$，并且
\[
\begin{aligned}
&\left.\frac{\partial}{\partial\epsilon_1}
L(\bar Q_{\bar g,Q_W}(\epsilon),Q_W)(o)\right|_{\epsilon=0}\\
&\qquad=H^{(1)}(\bar g)(o)[y-\bar Q(w)]=D_Y^{(1)*}(P)(o),\\
&\left.\frac{\partial}{\partial\epsilon_2}
L(\bar Q_{\bar g,Q_W}(\epsilon),Q_W)(o)\right|_{\epsilon=0}\\
&\qquad=\bar H_n^{(2)}(\bar g,Q_W)(o)[y-\bar Q(w)]=\bar D_n^{(2)}(P)(o).
\end{aligned}
\]
因此，给定 $(\bar Q_0,\bar g_0)$ 的初始估计 $(\bar Q_{n,0},\bar g_n)$，让 $\bar Q$ 遍及由
\[
\bar Q_{n,0}(\epsilon):=\expit\left[\logit(\bar Q_{n,0})
+\epsilon_1H^{(1)}(\bar g_n)+\epsilon_2\bar H_n^{(2)}(\bar g_n,Q_{W,n})\right]
\]
及 $-\infty<\epsilon_1,\epsilon_2<+\infty$ 确定的参数子模型，选择经验风险 $P_nL_1(\bar Q)$ 的极小点，即可得到 $\bar Q_0$ 的第一次更新估计 $\bar Q_{n,1}$。$\epsilon_1$ 和 $\epsilon_2$ 的最优值，可以很容易地作为logistic回归拟合中的回归系数估计而获得。这一回归的结局为 $Y_i$，协变量为 $H^{(1)}(\bar g_n)(O_i)=A_i/\bar g_n(W_i)$ 以及
\[
\bar H_n^{(2)}(\bar g_n,Q_{W,n})(O_i)=\frac{2A_i}{\bar g_n(W_i)}
\left[1-\frac{\bar g_{n,\mathrm{NP}}(W_i)}{\bar g_n(W_i)}\right],
\]
偏移项（offset）为 $\logit(\bar Q_{n,0}(W_i))$，并且只使用满足 $A_i=1$ 的数据点。不难看出，如果再考虑经过 $\bar Q_{n,1}$ 的二阶最不利子模型，试图进一步更新它，将不会产生任何变化。因此，此时算法在一步后即终止，故 $\bar Q_n^*=\bar Q_{n,1}$，式 (26.8) 必须成立。最终，$\psi_0$ 的2-TMLE为
\[
\Psi(\bar Q_n^*,Q_{W,n})=\frac1n\sum_{i=1}^n\bar Q_n^*(W_i).
\]
与前述1-TMLE相比，这一估计量带来了明确的理论改进；下面报告的初步模拟结果表明，这些改进确实很好地转化为了实际收益。

在上述算法中，完全不需要更新 $Q_{W,0}$ 的估计。这是因为我们选择非参数极大似然估计量（NPMLE）$Q_{W,n}$ 作为初始估计量，并对 $Q_W$ 使用对数似然损失。如果这两个选择中的任何一个不同，通常就必须使用适当扰动子模型，迭代更新 $Q_{W,0}$ 的估计，使所得靶向估计 $Q_{W,n}^*$ 满足
\[
P_nD_W^{(1)*}(\bar Q_n^*,Q_{W,n}^*)=0,
\]
从而得到 $\psi_0$ 的2-TMLE $\sum_w\bar Q_n^*(w)Q_{W,n}^*(w)$。
